\clearpage
\begingroup
\fancyhead[L]{\small THEORIE DER BEWEGUNG DES MONDES.}
\renewcommand{\headrulewidth}{0pt}
\par\noindent\hspace*{\parindent}Wir nennen die Grössen \(n\), \(e\), \(\varepsilon\), \(g\), wenn sie in bestimmte Zahlen
multiplicirt sind, Grössen von Einer Dimension; Producte aus zweien solchen
Grössen, Grössen von zwei Dimensionen; Producte aus dreien, Grössen von
drei Dimensionen u.\,s.\,w. Bei der Formel für die Länge, Breite, \(p\) u.\,s.\,f.
nennen wir alle Grössen von Einer Dimension, von zweien, dreien u.\,s.\,w.
Theile der ersten, zweiten, dritten Ordnung u.\,s.\,f. Hier ist sogleich klar,
dass, um die Differentialgleichung des 2\textsuperscript{ten} Grades für \(\theta\) bis auf die Grössen
der 5\textsuperscript{ten} Ordnung incl. genau zu haben, man die Werthe von
\(\dfrac{B^3}{R^3}\), \(\dfrac{\mathrm{d}u^2}{\mathrm{d}v^2}\) bis
auf die zweite Ordnung, den Werth von \(\theta\) aber bis auf die 3\textsuperscript{te} genau haben
müsse. Wir haben aber nur die Theile, deren Argumente \(\lambda\) und \(2E-\lambda\) sind,
angewandt, die resp. von der ersten und zweiten Ordnung sind; es ist also
nothwendig, um jene Differentialgleichung bis auf die 5\textsuperscript{te} Ordnung genau zu
haben, dass wir auch auf die 5 aus der Integration so eben entstandenen
Theile der 3\textsuperscript{ten} Ordnung Rücksicht nehmen. Auf diese Weise wird noch hinzugefügt
\begingroup\small
\setlength{\arraycolsep}{5pt}
\renewcommand{\arraystretch}{1.2}
\[
\begin{array}{l@{\qquad}l}
\multicolumn{1}{c}{\text{vermöge des Theils}}&
\multicolumn{1}{c}{\text{[zur Differentialgleichung}}\\
\multicolumn{1}{c}{\text{[von }\theta\text{:]}}&
\multicolumn{1}{c}{\text{für }\theta\text{:]}}\\
&\multicolumn{1}{c}{\tfrac32nng\times}\\
+\tfrac58eeg\sin(2M-\lambda)&-\tfrac58ee\sin(2E-2M+\lambda)\\
-\tfrac98n\varepsilon g\sin(\lambda+a)&+\tfrac98n\varepsilon\sin(2E-\lambda-a)\\
+\tfrac98n\varepsilon g\sin(\lambda-a)&-\tfrac98n\varepsilon\sin(2E-\lambda+a)\\
+\tfrac38n\varepsilon g\sin(2E-\lambda+a)&-\tfrac38n\varepsilon\sin(\lambda-a)\\
-\tfrac78n\varepsilon g\sin(2E-\lambda-a)&+\tfrac78n\varepsilon\sin(\lambda+a).
\end{array}
\]
\endgroup
\noindent Folglich kommt zu dem Werthe von \(\theta\) noch hinzu:
\[
+g\left\{\tfrac{15}{64}nee\sin(2E-2M+\lambda)\right\}
\]
\SourceLeaf{637}\VIIInlineSourceHead{}{THEORIE DER BEWEGUNG DES MONDES.}{631}
\noindent und es werden die Coefficienten von
\[
\begin{array}{r|l}
\multicolumn{1}{c}{g\times\sin}&\\
\lambda+a&-\tfrac98n\varepsilon\left(1-\tfrac{23}{24}n\right)\\
\lambda-a&+\tfrac98n\varepsilon\left(1-\tfrac18n\right)\\
2E-\lambda+a&+\tfrac38n\varepsilon\left(1+\tfrac{19}{8}n\right)\\
2E-\lambda-a&-\tfrac78n\varepsilon\left(1+\tfrac{65}{56}n\right).
\end{array}
\]
\par\noindent\hspace*{\parindent}Um diese Formel mit den [\textsc{Mayer}schen] Tafeln vergleichen zu können,
sind erst noch verschiedene Reductionen damit vorzunehmen. —

\par\noindent\hspace*{\parindent}I. Müssen wir darin statt \(E\), d.\,i. der Entfernung des Mondes von einer
fingirten mit einer mittlern Geschwindigkeit, die der des Mondes proportional
ist, laufenden Sonne, \(E^{\prime}\), die Entfernung des Mondes von der wahren Sonne
einführen.

\par\noindent\hspace*{\parindent}II. Anstatt \(\lambda\), der Entfernung des Mondes in der Ekliptik von einem
fingirten und mit einer der des Mondes proportionalen Geschwindigkeit
laufenden Knoten, haben wir \(\lambda^{\prime}\), die Entfernung des Mondes in der Bahn
von einem gleichförmig laufenden Knoten in Rechnung zu bringen.

\par\noindent\hspace*{\parindent}III. Gleichfalls statt \(M\) ist \(M^{\prime}\), Entfernung vom gleichförmig vorrückenden
Apogäum und

\par\noindent\hspace*{\parindent}IV. Statt \(a\) ist \(a^{\prime}\), die wirkliche mittlere Anomalie der Sonne zu setzen.

\par\noindent\hspace*{\parindent}Endlich V. ist aus \(\theta\), der Tangente der Breite, die Breite selbst, \(\beta\),
abzuleiten. Zu diesen Endzwecken haben wir folgende Gleichungen:
\[
\begin{aligned}
\text{I.}\quad E={}&E^{\prime}-2\varepsilon\sin a+\tfrac54\varepsilon\varepsilon\sin2a
 +2ne\sin M+\tfrac34nee\sin2M\\
&-\tfrac{11}{8}n^3\sin2E+\tfrac{15}{4}nne\sin(2E-M)
 +\tfrac14ngg\sin2\lambda-3nn\varepsilon\sin a
\end{aligned}
\]
\[
\begin{aligned}
\text{II.}\quad\lambda={}&\lambda^{\prime}
 -\left(\tfrac14gg-\tfrac18g^4\right)\sin2\lambda^{\prime}
 +\tfrac{1}{32}g^4\sin4\lambda\\
&-\zeta\bigl\{2e\sin M+\tfrac34ee\sin2M-\tfrac{11}{8}nn\sin2E\\
&\qquad+\tfrac{15}{4}ne\sin(2E-M)+\tfrac14gg\sin2\lambda-3n\varepsilon\sin a\bigr\}
\end{aligned}
\]
\[
\text{III.}\quad M=M^{\prime}+2e\eta\sin M
\]
\[
\text{IV.}\quad a=a^{\prime}+2ne\sin M
\]
\[
\text{V.}\quad\beta=\theta-\tfrac13\theta^3+\tfrac15\theta^5.
\]
\clearpage
\endgroup
